\chapter{Network Analysis}
\label{chap:networkAnalysis}

This chapter examines the network traffic generated by the Ray-Ban Meta glasses and the Meta View app. Chapter \ref{chap:methodologyExperiment} describes the traffic analysis methods and the testbed for capturing the traffic, which includes both active usage scenarios of the glasses (touch-based commands, voice-based commands, and pairing) and idle traffic. It also explains the different types of parties in the analysis. The focus is to assess the privacy implications of network traffic captured during different interactions, including user profiling, location inference, and bystander exposure. Specifically, to identify what patterns can be observed, determine whether any personal or bystander information is exposed, and pinpoint the potential privacy and security risks. Due to most traffic being encrypted, the analysis focuses on metadata, such as domain names, packet sizes, destinations, and protocols, to infer such information.
 

\section{General Network Analysis}
The captured network traffic from interactions with the glasses reveals several potential concerns related to user data and privacy worth mentioning. For example, a general analysis across all experiments showed that an obvious pattern can be found in the traffic volume and protocols used for media transfer, which will be detailed later. This section examines three prevalent trends observed across the entire data collection. The first is protocol usage, which analyzes the types of protocols captured in the experiments. The second trend explores the general flow of traffic, including the size distribution of network packets and the sequence of packets in the traffic. The last aspect focuses on the destination of the packets.

\subsection*{Protocols}
\label{subsec:protocol}
Figure \ref{fig:protocol} shows all protocols captured in the active interactions experiments and their readability with examples based on the payload. A protocol type is considered unreadable if the packet's content is not human-readable. A packet is partially readable if the payload consists of readable elements, usually the domain names, with the rest being unreadable. Readable is when the majority of the payload is human-readable. 

Packets using the User Datagram Protocol (UDP) are an example of partially readable traffic. Within the payload of some UDP packets, a readable string was found, SIGPATHSREPCERTINDX. A dissected interpretation of this string could be SIG(nature) PATHS REP CERT(ificate) IN(de)X, which indicates that the packet was involved as part of a certificate exchange. However, the captured packets, as seen in the figure, predominately relied on Transmission Control Protocol (TCP), Transport Layer Security (TLS) versions 1.2 and 1.3, and Quick UDP Internet Connections (QUIC). Their payloads are typically encrypted, making them unreadable in raw form. Interestingly, there was an observable difference in the frequency of TLS handshakes between the experiments for touch-based and voice-based commands. In touch-based commands, TLS handshakes were consistent throughout the capture, whereas in voice-based commands, they only really appeared in the first two minutes of the capture or the last two minutes.

\begin{figure}[H]
    \centering
    \includegraphics[width=\textwidth, trim=0 390 0 70, clip]{Figures/NV_protocols.pdf}
    \caption{First-party traffic volume categorized based on the packet's protocol and readability. The accompanying Wireshark snapshots show example protocols for each type of readability, along with what can be read in their payloads.}
    \label{fig:protocol}
\end{figure}
 
Nevertheless, the accompanying metadata can still lead to information exposure. For example, in the information column, the Server Name Indication (SNI) fields in TLS `Client Hello' messages revealed repeated connections to domains such as graph.facebook.com as well as Spotify and WhatsApp APIs. Similarly, recurring destination domain names, including graph.facebook.com and spclient.wg.spotify.com, are visible in Domain Name System (DNS) queries. Despite the encryption of the payload, repeated connections to particular domains, especially under a certain time interval, can disclose details about user activities and application usage. 

Additional protocols, such as Address Resolution Protocol (ARP) and Dynamic Host Configuration Protocol (DHCP), revealed device identifiers and MAC and IP addresses. More specifically, from ARP, the IP addresses of the phone and network were exposed, and the MAC addresses of the phone, network, and glasses were exposed. The exposed IP addresses of the network connection provide a general location of the user, and the MAC addresses of the devices were consistent throughout the experiments, which allows potential user fingerprinting. In all captures, the ALFA network adapter appeared as Alfa\_b2:e8:9a, showing how unique device identifiers appear regularly in network traffic. Another protocol, Dynamic Host Configuration Protocol (DHCP), used for assigning IP addresses, has DHCP requests packets that expose additional information, including the phone as android-dhcp-14 and the glasses as Ray Ban-s-S247. Such information could also aid in fingerprinting connected devices.

Furthermore, protocols such as Simple Service Discover Protocol (SSDP) and Multicast DNS (MDNS), which are for service or device discovery, are visible to all devices on the same network. As a result, local network information and the phone's discovery and connectivity attempts were on full display. In the captures, these protocols often had queries for the service name \_spotify-connect.\_tcp.local, which shows that the phone was running Spotify and tried to use this service name to find other devices that could play Spotify on the same network. This visibility of service-discovery traffic increases the risk of passive monitoring by exposing available services and device identities. 

\subsection*{Traffic Volume and Packet Size}

\begin{figure}[b]
    \centering
    \includegraphics[width=0.9\textwidth]{Figures/NV_packet_distribution.pdf}
    \caption{The packet distribution, which shows a comparison of the number of packets that are of a specific size for all traffic captured in the five rounds of voice-based commands experiments.}
    \label{fig:NV_packet_distribution}
\end{figure}

Another aspect of observation pertains to the overall traffic flow, where traffic peaks or fluctuations often correlate with specific application usage. For example, the largest spikes typically occur during media transfers, although experiment-specific patterns are discussed in detail in the subsequent sections. In order to have a better understanding of these patterns, the protocol-specific traffic volumes were analyzed and found to have recurring sequences that can also be tied to specific interactions. For instance, the interaction for media transfer typically involves four packets of Internet Control Message Protocol version 6 (ICMPv6) packets, followed by DNS queries to Facebook, which is then continued with TCP and TLS protocols with large packet sizes. Typically, most packets are relatively small, fewer than 200 bytes.

In contrast, media transfers frequently involve packet sizes above 1000 bytes, often around 1446 bytes. Figure \ref{fig:NV_packet_distribution} illustrates the packet size distribution for the voice-based command experiments. The two highest peaks in the figure are packets with sizes less than 200 and packets with sizes greater than 1400 bytes. Even without decryption, this difference in traffic volume allows an observer to infer user activity. It should be noted that these ICMPv6 packets appeared exclusively during media transfer, which simplifies the identification of this specific interaction, even if the actual content remains encrypted.

\subsection*{Traffic Destination}
Throughout the experiment, there were continuous interactions with public servers. First-party traffic dominates the traffic volume, communicating with Facebook APIs (e.g., graph.facebook.com). Furthermore,  there was frequent activity involving Spotify APIs, such as spclient.wg.spotify.com. While some Google APIs, such as android.googleapis.com, were detected, they appeared to be related either to Google Cloud services supporting Spotify or system services for the device. Many of the destinations of these addresses were located in Switzerland, allowing adversaries to get an idea of the user's general location.

It is worth noting that none of the identified third-party IP addresses were detected to be from prominent advertising or tracking domains. However, the Meta View app analysis in Chapter \ref{chap:appAnalysis} of the glasses' permission requests and the glasses' data collection policies \cite{GlassesPrivacyPolicy} indicate that Meta does collect user activity data. Therefore, the absence of explicit tracking does not eliminate the presence of tracking. However, this does imply that users have less control over their data, as they lack full transparency, making it difficult to know which activities are monitored, what data is being gathered for analytics, and how they may be used.


\section{Active Traffic}
Active traffic refers to the traffic generated during user interactions through touch-based commands, voice-based commands, and pairing sequences between the device and the companion app. The following subsections dive into each category of user interaction experiments, explaining unique traffic patterns and potential privacy risks for each experiment.

\subsection{User interactions - Touch commands}

The experiment for touch-based commands is primarily comprised of commands for Media Capture and Music Playback. Across all five repeated trials, packets sent to public destinations account for over 80\% of the total traffic volume. Among these public addresses, first-party traffic dominates the contribution to the traffic volume, as seen in Figure \ref{fig:NT_bar_log_combined_different_party_traffic}, exceeding 2.8 MB on average per trial and contributing more than 98\% of the public traffic volume. 

\begin{figure}[H]
    \centering
    \includegraphics[width=0.9\textwidth]{Figures/NT_bar_log_combined_different_party_traffic.pdf}
    \caption{The mean public network traffic captured across the five repeated trials of the touch-based commands experiments, with the data categorized by party type.}
    \label{fig:NT_bar_log_combined_different_party_traffic}
\end{figure} 

\begin{figure}[b]
    \centering
    \includegraphics[width=0.9\textwidth]{Figures/NT_line_all_traffic.pdf}
    \caption{Log line plot of public traffic captured during the five repeated trials of the touch-based commands experiments, with the data separated by trial.}
    \label{fig:NT_line_all_traffic}
\end{figure}

The bar graph in Figure \ref{fig:NT_bar_log_combined_different_party_traffic} shows that there are two prominent spikes in the traffic, which appear even more clearly in Figure \ref{fig:NT_line_all_traffic}. From the figure, the experiments have a relatively consistent pattern throughout the timeline, as indicated by the similar general behavior of the plotted lines. Due to the majority of traffic being first-party traffic, the patterns in the total traffic volume are analogous to that of first-party traffic and its volume trends. The first spike, as indicated with the magenta line, occurs during the initial power on of the glasses and opening of the Meta View app. The duration of this interaction typically spans from minute 0.5 to 1.5. The second prominent peak, as indicated by the cyan line, is during the transfer of the captured media from the glasses to the phone, which is expected as the media transport is through the network and involves bulk data transfer. In comparison, capturing media and playing music did not generate large traffic bursts. This striking difference in traffic volume demonstrates that even encrypted payloads could expose an approximate sense of user activity and that external entities could still determine certain user activities based on the traffic volume.

\begin{figure}[b]
    \centering
    \includegraphics[width=0.8\linewidth]{Figures/NT_line_support_party_traffic.pdf}
    \caption{Support-party traffic captured during the five repeated trials of the touch-based commands experiments, with the data separated by trial.}
    \label{fig:NT_line_support_party_traffic}
\end{figure}

On the other hand, support-party traffic was nearly negligible at under 5 KB on average per trial, which is less than 0.2\% of traffic volume. Figure \ref{fig:NT_line_support_party_traffic} shows that Trial 5 was somewhat inconsistent with the other trials. Upon inspection, it was noticed that there was an absence of Continuous Data SSL packets, which were present in all other trials. This inconsistency was likely due to network factors or interrupted connection, as the packets were highlighted in black and red, indicating failed connections and packet losses.

The destination of all traffic captured was a single WhatsApp service provider. These packets consisted mainly of TCP ACK packets and some SSL packets. This finding is unexpected, given that touch commands do not interact with nor involve WhatsApp. Looking at the figure, the interaction with support-party traffic appears consistently during the initial pairing, media capture, and when playing music on Spotify. None of these interactions require WhatsApp. As shown in Figure \ref{fig:NT_line_support_party_traffic}, this behavior was consistent across all trials, each with similar traffic volume. Focusing on the existence of this address during media capture, the figure suggests that the app, for some reason, initiates a connection to the WhatsApp API whenever the wearer takes a photo or video, even though such a connection would typically be expected only when attempting to upload or share the media on WhatsApp. According to Meta's Terms of Privacy \cite{MetaPrivacyPolicy}, images captured by the glasses or interactions made with the glasses are not stored unless uploaded onto one of their social media platforms or sent to their servers during the interaction. However, it is plausible that, by leveraging WhatsApp as an intermediary, these images could then be sent to their servers for analytics or other purposes without violating their terms of service, which would be a complete privacy invasion for users. This behavior was observed exclusively during physical interactions with the glasses and was absent during the voice-command experiments. A possible explanation could be Meta's policy, which states that they already collect the information from voice-enabled features, reducing the need for additional transmissions during voice interactions. This unexpected behavior raises questions about whether the app unnecessarily transmits data to servers during media capture without the user's awareness.


There was minimal third-party traffic, below 60 KB on average per trial, contributing to less than 2\% of the traffic volume. This indicates that there are only a few external dependencies when physically interacting with the glasses. As shown in Figure \ref{fig:NT_bar_log_combined_different_party_traffic}, minutes 3.5 and 4.5, corresponding to the start of music playback and change of the music track, had the most notable peaks in third-party traffic occurrence. Unsurprisingly, these surges involve communication with Spotify servers. During these time intervals, there was also nearly no first-party traffic, suggesting that music playback does not rely on Meta servers. 

In order to analyze the extent of information exposure, the traffic was categorized by protocol to assess the readability of the network traffic. This is illustrated in Figure \ref{fig:NT_bar_log_all_protocol_traffic}, which shows that most packets were not human-readable. Few of the partially readable or readable traffic included local protocols, such as DHCP, MDNS, and SSDP. 

\begin{figure}[H]
    \centering
    \includegraphics[width=0.9\textwidth]{Figures/NT_bar_log_all_protocol_traffic.pdf}
    \caption{All network traffic captured across the five repeated trials of the touch-based commands experiments, with the data categorized by readability and protocol.}
    \label{fig:NT_bar_log_all_protocol_traffic}
\end{figure}

\newpage
Protocols for local networks can still reveal private details. MDNS and DHCP protocols revealed information about the services and devices in the network. For example, \_spotify-connect.\_tcp.local indicates that Spotify Connect is being queried, allowing direct music streaming from Spotify. This protocol only reveals the IP address of the phone trying to find other devices compatible with running Spotify. However, this still informs third parties that the phone is running Spotify. There is one protocol, SSDP, with readable payloads. They are packets with (Hypertext Transfer Protocol) HTTP content. The readable content includes MAN: "ssdp:discover," specifying it as a discovery message. The type of service being searched is also visible from search target (ST), ST: urn:dial-multiscreen-org:service:dial:1. Discovery and launch (DIAL) is a protocol used by devices to allow apps to be discovered and launched, meaning this packet is a discovery query. Both MDNS and SSDP packets in the captures are in two-minute intervals and can expose information about the types of devices and services available on the local network to unauthorized devices. 

\begin{figure}[b]
    \centering
    \includegraphics[width=0.9\textwidth]{Figures/NT_bar_log_third_party_protocol_traffic.pdf}
    \caption{Third-party traffic captured across the five repeated trials of the touch-based commands experiments, with the data categorized by readability and protocol.}
    \label{fig:NT_bar_log_third_party_protocol_traffic}
\end{figure}

The rest of the protocols handle public traffic. DNS protocols are one of the two partially readable protocols. The queries revealed the hostnames of the targeted domains. There were frequent queries for Facebook domains, such as graph.facebook.com or gateway.facebook.come. The consistent DNS throughout the capture not only reveals which application is active but can also indicate server communication patterns that could be used to infer user behavior. The other partially readable public protocol is UDP. The UDP packets had a payload that started with the tag SpotUdp0, which appears to be an application-specific identifier, which allow observers to infer application usage. 

There was an interesting singular packet with protocol TLSv1 was directed to gateway.facebook.com. This was in Trial 1 around minute 5.0, near the capture's end. Although no sensitive user data is revealed, metadata, including the hostname, indicated in the SNI or the cryptographic details, such as the supported cryptographic algorithms or TLS versions, were clearly revealed. However, its presence in the capture raises questions about potential server configurations that would result from the deprecated version of the protocol.

Focusing on third-party traffic, only five protocols were identified: NTP, QUIC, TCP, TLSv1.2, and TLSv1.3. These protocols are visualized in Figure \ref{fig:NT_bar_log_third_party_protocol_traffic}. Most of the traffic, as illustrated in the figure, is unreadable, reflecting the substantial encryption in these communications. The partially readable traffic in TLSv1.3 only consisted of `Client Hello' messages that revealed the hostnames of the phone that wanted to establish a secure connection. More than half of these hostnames were directed to Spotify, such as clienttoken.spotify.com and login5.spotify.com, which are related to authentication tokens or the login procedure for Spotify. The rest of the partially readable traffic was mainly related to Google, which had queries to hostnames such as www.google.com. As a result, although the exact content is unknown, the payload still exposed the references to the domains. Spotify and Google both had packets with protocols TCP, TLSv1.2, and TLSv1.3. Interestingly, Google was the sole contributor to QUIC packets.

\begin{figure}[b]
    \centering
    \includegraphics[width=\textwidth]{Figures/NT_line_combined_keeping_outlier_third_party_traffic.pdf}
    \caption{Third-party traffic that contributed more than 300 bytes captured across the five repeated trials of the touch-based commands experiments, with the data categorized by IP addresses. The legend displays each IP address along with its corresponding ISP or Service if known, and the Total Number of Bytes.}
    \label{fig:NT_line_combined_keeping_outlier_third_party_traffic}
\end{figure}

To better illustrate the third-party traffic destinations, Figure \ref{fig:NT_line_combined_keeping_outlier_third_party_traffic} shows the total traffic volume for third parties for different IP addresses that contributed over 200 bytes. While many of the addresses are directed to Google, Spotify produced the largest amount of third-party traffic in this experiment. Spotify had more than 40 KB on average per trial, meaning that it accounted for more than 70\% of all third-party traffic. The reason is that the average traffic volume of Spotify addresses was much higher than that of Google. It is noted that although the IP addresses of the traffic for the orange lines are registered to Google, the traffic is directed to Spotify services, as Spotify is using Google Cloud to host its services. This concept was explained in \ref{subsec:networkTraffic}. The orange lines in Figure \ref{fig:NT_line_combined_keeping_outlier_third_party_traffic} show that Spotify was more active during the last two minutes of the experiment during music playback. Google was the next largest third-party destination; however, the packets were either related to Google Play services or were there for connectivity. The presence of each of those IP addresses was also inconsistent throughout the trials, which likely indicates that their existence was not from the app. For example, the blue spike belonging to Google at minute 3.0 was solely from Trial 3; no other trials had this specific IP address and did not display such a spike at that minute. The protocol for this Google address was QUIC, and, as illustrated in Figure \ref{fig:NT_bar_log_third_party_protocol_traffic}, all QUIC payloads are unreadable. It is notable that no destination addresses were linked to any trackers or advertisers, though this does not exclude their existence; only none were found in the domains and readable content. 
    
The geological location of the third-party traffic directed towards Spotify is in CA, US, as seen in Figure \ref{fig:NT_maps}. The traffic to Google is also sent to the US, although Google does have some addresses in Switzerland. There were numerous destination addresses that pointed to Switzerland, the same country where the experiment took place. Although these IPs were minor in terms of traffic volume, the location of these addresses can still be used by an adversary to determine the approximate geological location of the user. Third-party entities can also identify the services or applications the user uses based on the IP addresses. Unexpectedly, there is a destination in the figure shown to be in Taiwan, Taipei, which is from Trial 3. The IP address of this destination belongs to Google Cloud services. There is no data payload in the four packets captured; they are purely acknowledgment packets. There is another notable cloud service in the US, where the address is registered to Amazon. The few packets to this address were only found in Trial 2 and were all acknowledgment packets as well.

\begin{figure}[H]
    \centering
    \includegraphics[width=\textwidth, trim=0 330 0 270, clip]{Figures/NT_maps.pdf}
    \caption{Destination addresses of all traffic captured across the five repeated trials of the touch-based commands experiments.}
    \label{fig:NT_maps}
\end{figure}

An additional unique destination was found in Ireland. This was from a singular packet that is part of a TLS handshake from Trial 1. Specifically, it was a TLS server certificate message. Its purpose is to prove the server's identity to the client to establish trust and contain the public key. Unlike others, it was partially readable. In the packet, the root CA in the certificate chain was identified as Amazon Root CA, meaning it is managed by Amazon. There was another CA, Starfield Services Certificate Authority, from Starfield Technologies, Inc., which is most likely an intermediate CA in the certificate chain that signed the server's certificate. References to certificate authority (CA) endpoints could be identified, including http://ocsop.rootg2.amazontrust.com/rootg2.cert0= and http://crl.rootg2.amazontrust.com/rootg2.crl0, which are link points to the certificate of and Certificate Revocation List (CRL) for the root CA, Amazon Root CA. The location where the CA is registered was also shown as Scottsdale, Arizona, US. Third-party entities can infer the server's identity from this packet using the IP address from the metadata, ec2-54-194-107-68.eu-west-1.compute.amazonaws.com. Based on the domain, this server appears to be part of Amazon's computing services in the West EU. 
Based on the Client Hello in the Information column of the packet, sdk.push.samsung.com, this packet is most likely related to push notifications for the Samsung phone. Nevertheless, attackers could use the information from the packing, including the server's hostname, the IP address, and certificate details, which include links to CA validation endpoints and information about the root and intermediate CAs, to map the validation infrastructure used by the CA, providing attackers the ability to target certificate validation servers.


\subsection{User interactions - Voice commands}

\begin{figure}[b]
    \centering
    \includegraphics[width=0.9\textwidth]{Figures/NV_bar_log_combined_different_party_traffic.pdf}
    \caption{The mean public network traffic captured across the five repeated trials of the voice-based commands experiments, with the data categorized by party type.}
    \label{fig:NV_bar_log_combined_different_party_traffic}
\end{figure}

From the combined capture of all five repeated trials of the experiment, nearly 70\% of the traffic volume for the voice-based commands experiment was directed to public destinations. As illustrated in Figure \ref{fig:NV_bar_log_combined_different_party_traffic}, most of this public traffic is first-party traffic, with an average of more than 3.5 MB on average per trial and comprising over 95\% of the total traffic. Although support-party traffic rose in terms of volume, with 50 KB on average per trial, which is nearly 7.5 times the volume from touch-based commands, it still contributes to less than 1\% of the public traffic volume. The figure shows a slight increase in support-party traffic during the last 2 minutes, though it is mostly around minute 10. This high frequency of traffic reflects the voice command that sends a voice message to a person. There was also an increase in third-party traffic, averaging 120 KB per trial. However, considering that the duration of the experiment is twice as long as touch-based commands and the average traffic volume also doubles the amount of touch-based commands, no extra third-party traffic is being generated from the different command types. In essence, the similar proportion of third-party traffic volume suggests that the traffic is driven more by user interaction than by the specific input method.


The voice command experiments execute a command every thirty seconds, starting from minute 0.5, which is reflected in Figure \ref{fig:NV_scatter_log_public_traffic}. The scatter plot displays the traffic volume aggregated every ten seconds, with the cyan lines indicating the thirty-second intervals. Most dots with higher traffic volume are located on the cyan lines. This means that an adversary could presume that when the traffic volume is roughly $10^5$, it is likely that the user is executing a voice command on the glasses. For many of the voice commands, the data points across different trials are relatively close, which suggests a higher reliability for the adversary to infer that traffic volume corresponds to user interaction. This pattern was unique to voice-based commands and was not observed in touch-based commands. This adds to the ability of the adversary to determine which method of interaction was used, voice-based or touch-based.

\begin{figure}[H]
    \centering
    \includegraphics[width=0.9\textwidth]{Figures/NV_scatter_log_public_traffic.pdf}
    \caption{Public network traffic captured across the five repeated trials of the voice-based commands experiments, with the data grouped in ten-second intervals.}
    \label{fig:NV_scatter_log_public_traffic}
\end{figure}

\begin{figure}[b]
    \centering
    \begin{minipage}{0.49\textwidth}
        \includegraphics[width=\textwidth]{Figures/N_quic_send_message.png}
    \end{minipage}
    \begin{minipage}{0.48\textwidth}
        \includegraphics[width=\textwidth]{Figures/N_quic_date.png}
    \end{minipage}
    \caption{Wireshark capture of QUIC packets when asking to send a message vs. date in Trial 5 of the voice-command experiment.}
    \label{fig:quic_comparison}
\end{figure}

A more specific traffic pattern noted during the experiments was based on the protocol used. There is a visible pattern from the QUIC protocol between the phone and the Meta servers. There are consistent, consecutive 5 to 7 packets, but usually 6 packets, with the QUIC protocol, which contains a payload that is transmitted immediately when a command is executed. This sequence of traffic can be easily identified because, in idle capture, there were nearly no consecutive 6 QUIC packets across the 8-hour period. This identifiable traffic pattern allows for better tracking of user interactions with the glasses. It is also possible to have an idea of what command the user provided. Figure \ref{fig:quic_comparison} shows a side-by-side comparison of two example commands in Wireshark from Trial 5, one asking for the AI to send a message and the other asking for the date. The images in the figure only include packets with the protocol QUIC, and the two columns displayed are the time of the packet capture and the information of the packets. The command to ask the AI to send a message is at minute 9.5 (570 seconds), as shown in the image on the left, while asking for the date is at minute 2.5 (150 seconds), which is the image on the right.

As seen in the right of Figure \ref{fig:quic_comparison}, asking for the date usually results in the 6 consecutive QUIC packets that contain payload followed by 4 empty QUIC packets, which is then continued with nearly 20 QUIC packets with payload. The sequence is then separable from when asking Meta to send a message to a person, as seen in the left of Figure \ref{fig:quic_comparison}, which is due to the back and forths of the glasses needing to know who the receiver of the message is, the content of the message, and the final confirmation to send. There would be 6 to 8 consecutive QUIC packets with payload followed by 4 without, then once again 6 to 8 with and 4 without, forming a cycle until the interaction is complete. This sequence is then also different in comparison to asking for the weather, which usually results in the 6 consecutive QUIC packets containing payload followed by 3 to 4 empty QUIC packets, which are continued with 8 or 9 consecutive QUIC with payload. These are general patterns, so there could be a couple of empty QUIC packets within the group of QUIC packets with payload. However, by observing the number of packets with specific protocols transmitted, a third party could roughly estimate the user's interaction. However, this depends on the interaction as multiple other commands have similar traffic sequences as when asking for the date, whereas asking for the weather and asking to send a message result in more unique traffic patterns.


Figure \ref{fig:NV_line_log_all_traffic} provides a better overview of the general traffic flow across different trials, illustrating the timeline of all public traffic. The figure shows that the trials were rather consistent with one another, with the only discrepancy during minute 9.0, where Trial 4, the red line, and Trial 5, the blue line, experienced a large increase, while Trial 3, the green line, had a drop in traffic volume. This is also reflected in the average traffic volume for each trial, with a median of 4 MB per trial for trials without the spike, 8 MB for Trial 4, and 6 MB for Trial 5. The spikes are mainly due to first-party traffic from the glasses to the phone during media import, as seen in the surge in minute 9.0. There was no observable reason as to why the trials had such differences in volume during this period. There are two time periods in the figure with a noticeable surge in traffic volume. The first, smaller increase between minutes 0.5 and 2, is the interaction to perform the initial powering on the glasses and opening of the Meta View app. This increase was consistent throughout the trials. The second larger but less consistent increase is at minute 9.0, which coincides with the transfer of media captures from the glasses to the phone. 

\begin{figure}[H]
    \centering
    \includegraphics[width=0.9\textwidth]{Figures/NV_line_log_all_traffic.pdf}
    \caption{Public network traffic captured during the five repeated trials of the voice-based commands experiments, with the data separated by trial.}
    \label{fig:NV_line_log_all_traffic}
\end{figure}

Diving into the captured support-party traffic, one IP address pointed to Instagram, and two were to WhatsApp. WhatsApp was relatively more consistent throughout the trials, while Instagram's IP address was only in three of the five trials. Both of the traffic destinations were largely unreadable. However, the same packet with the same IP address was present for each trial, with the Info column Client Hello (SNI=media-zrh1-1.cdn.whatsapp.net). This specific packet was only captured once during each trial at the exact time of minute 10.33. The interaction during this time was when receiving a message on WhatsApp. Throughout the trials, there were no other packets with this exact Info column or IP address other than during this interaction. As a result, a malicious person would be able to precisely identify when a person received a message on WhatsApp and, if able to gain physical access to the glasses (with the phone lock feature off), be able to obtain the message sent to the user.

Figure \ref{fig:NV_line_log_third_party_traffic} illustrates the third-party traffic volume for different trials. The traffic between the experiment trials aligns surprisingly well with one another, with only some inconsistencies, suggesting that the traffic generated is predictable and closely tied to the commands being executed. As a result, it can be used by attackers to infer the type of interactions the user performs. The first inconsistency is during the initial powering on and opening of the Meta View app, which is between minutes 0.5 and 1.5. The other inconsistency is that in Trial 3, where a much larger spike is observed during the command to make a video recording in minute 8. However, it seems that Trial 3 was an outlier for that interaction. The specific spike in Trial 3 is roughly 200 KB. However, of those, they are all unreadable traffic, with QUIC having roughly 160KB, followed by TCP with a little more than 30KB and TLS with 13KB bytes, as seen in Figure \ref{fig:NV_bar_log_third_party_protocol_traffic}, which shows the third-party traffic grouped by protocol and readability. The time between minutes 3.5 and 6.5 has the closest alignment of the trials, which starts at the beginning of music playback and ends during the timer interaction. The consistency in the traffic volume can be seen as a positive indicator that the traffic being generated for Spotify is most likely essential communication and is considered necessary for the app's functionalities, although analytics or advertisers could still be included. 

\begin{figure}[H]
    \centering
    \includegraphics[width=0.9\textwidth]{Figures/NV_line_log_third_party_traffic.pdf}
    \caption{Third-party traffic captured during the five repeated trials of the voice-based commands experiments, with the data separated by trial.}
    \label{fig:NV_line_log_third_party_traffic}
\end{figure}


The majority of all traffic captured was encrypted and, hence, not human-readable. This is the same for third-party traffic volume, as displayed in Figure \ref{fig:NV_bar_log_third_party_protocol_traffic}. The partially readable traffic from the protocol TSLv1.3 is roughly the same as in touch-based traffic, where it is `Client Hello messages,' which reveal the domain name the client is requesting a connection to. 

Besides the five protocols in touch-based commands, the figure also included some HTTP traffic. This HTTP traffic captured was not present in the experiment for touch-based commands. Out of the five trials conducted, this traffic was present in three of them. Each trial only had two HTTP packets, an HTTP GET request and an HTTP response, at different time periods. The phone sends a GET request to /generate\_204 HTTP/1.1, which is an endpoint for connectivity checks. The packet completely reveals the software making the request, in this case, Mozilla/5.0 (X11; Linux x86\_64) AppleWebKit/537.36 (KHTML, like Gecko) Chrome/60.0.3112.32 Safari/537.36, which shows that Chrome on a Linux device is making the request. The server responds with HTTP/1.1 204 No Content, meaning that network connectivity is confirmed, but no content is returned in the response. Nevertheless, the exact date and time are included in the response. The purpose of this traffic is solely for the device to do a connectivity check. Although no sensitive information is exposed, the software details of the device and the exact timestamps in the packet can still contribute to user profiling.

\begin{figure}[H]
    \centering
    \includegraphics[width=0.9\textwidth]{Figures/NV_bar_log_third_party_protocol_traffic.pdf}
    \caption{Third-party traffic captured across the five repeated trials of the voice-based commands experiments, with the data categorized by readability and protocol.}
    \label{fig:NV_bar_log_third_party_protocol_traffic}
\end{figure}

The third-party traffic volume still had the same main contributor, Spotify, as in touch-based commands, with Google being the second largest destination. Figure \ref{fig:NV_line_combined_keeping_outlier_third_party_traffic} shows the timeline of third-party traffic based on the IP address. This figure shows all third-party traffic that contributed more than 300 bytes of traffic. It shows that the main third-party destination belongs to Spotify, which has almost 70 KB per trial, nearly 60\% of third-party traffic volume, while Google had an average of 50 KB per trial. The traffic for Spotify occurred mostly during Music Playback, as expected, though it was consistently present throughout the whole experiment. Once again, the peak appears to be an anomaly in Trial 3. As mentioned, the minimal third-party interactions show that the app avoids unnecessary connections to other third parties, which is favorable for reducing attack vectors and protecting user data. However, the predictable reliance on Google and Spotify could be combined with other traffic patterns tied to user interactions, which could help attackers in traffic fingerprinting. 

\begin{figure}[H]
    \centering
    \includegraphics[width=\textwidth]{Figures/NV_line_combined_keeping_outlier_third_party_traffic.pdf}
    \caption{Third-party traffic that contributed more than 300 bytes captured across the five repeated trials of the voice-based commands experiments, with the data categorized by IP addresses. The legend displays each IP address along with its corresponding ISP or Service if known, and the Total Number of Bytes}
    \label{fig:NV_line_combined_keeping_outlier_third_party_traffic}
\end{figure}

\begin{figure}[b]
    \centering
    \includegraphics[width=\textwidth, trim=0 300 0 250, clip]{Figures/NV_maps.pdf}
    \caption{Destination addresses for all traffic captured across the five repeated trials of the voice-based commands experiments.}
    \label{fig:NV_maps}
\end{figure}

The geological location of most third-party traffic volume is routed to U.S.-based servers, as displayed in Figure \ref{fig:NV_maps}. It is interesting that despite a Google server being located in Switzerland, most of the traffic to Google IP addresses is routed to the US, which shows that Google's data centers are primarily located and operating in the US. During this experiment, most third-party addresses in Switzerland, illustrated in Figure \ref{fig:NV_maps}, were third parties that contributed negligibly in terms of traffic volume. However, its location still exposes the general location of the user.

An interesting observation about the destination of first-party traffic is that the IP addresses of these traffic point to locations in Germany and Italy, compared to the destination for touch-based commands being in Switzerland. A deeper analysis of the packets of these IP addresses show that they are part of a process called STUN/TURN, which helps devices connect to each other over the internet behind firewalls or routers. The TURN server for those IP addresses acts as a middleman when direct connections between the devices are not possible. 

In this case, Meta is routing the traffic through its TURN servers in Germany and Italy. As a result, most first-party locations are directed to Meta servers in Europe, more specifically Germany and Italy. The additional addresses only occur during the last minute of the experiment when there is an incoming phone call. This means that for most user interactions, the servers in Switzerland are able to handle the commands. However, phone calls require a more distributed system to handle communication more efficiently, causing the need for servers in Germany and Italy. As a result, based on the destination servers of first-party traffic, a third-party entity could also infer when the user is on a phone call. The map also shows that the location for support-party traffic and most third-party traffic destined for Zürich, Switzerland, which allows for an approximation of the user's location. 

\subsection{Pairing Sequence}

From initial hands-on usage of the glasses, it is revealed that the device would not allow pairing without an internet connection. The pairing process starts with the initial discovery and establishing a connection between the Ray-Ban Meta and the phone. Then, data is exchanged to authenticate and validate the devices through secured channels. The reason for this authentication is to prevent two devices from connecting to the glasses at the same time. Due to the fact that the glasses store images locally, if another person were to be directly able to connect to the glasses, they would be able to see those images after pairing them with the glasses, meaning a breach of user privacy. 

\begin{figure}[t]
    \centering
    \includegraphics[width=\textwidth]{Figures/NP_line_all_ip_address_traffic.pdf}
    \caption{Public network traffic that contributed more than 100 bytes captured in the initial paring experiment, with the data categorized by IP addresses. The legend displays each IP address along with its corresponding ISP or Service if known, and the Total Number of Bytes.}
    \label{fig:NP_line_all_ip_address_traffic}
\end{figure}

From the traffic capture, the split between public and private traffic volume was relatively even, with public network traffic representing more than 47\% of the total traffic volume, an estimate of 1.25 MB. Most public traffic volume stems from first-party traffic, contributing more than 98\% of the traffic volume, nearly 1.25 MB on its own. No support-party traffic was captured at all, and only about 20 KB of third-party traffic was recorded. 
Some communication with external cloud servers, mainly Google, was observed in the traffic. Figure \ref{fig:NP_line_all_ip_address_traffic} shows the drastic difference in traffic volume of the teal and green IP addresses compared to the rest of the addresses in the figure. Each cyan line in the plot indicates a new round of pairing attempts. From the initial overview, the traffic appeared to have no relative pattern due to the inconsistencies and differences in traffic volume. However, a deeper look into specific packets shows a certain sequence of traffic that can be observed during pairing between the devices. 

\begin{figure}[b]
    \centering
    \includegraphics[width=0.9\textwidth]{Figures/NP_bar_log_third_party_protocol_traffic.pdf}
    \caption{Third-party traffic captured during the initial pairing experiment, with the data categorized by readability and protocol.}
    \label{fig:NP_bar_log_third_party_protocol_traffic}
\end{figure}

It begins with a sequence of DNS queries. DNS lookups involved include ar.graph.meta.com and graph.facebook.com. these domains were also present in the captures of the experiments for touch-based and voice-based commands during the initial powering on of the glasses. This is then followed by a sequence of TCP handshakes, including packets with SYN or SYN-ACK, which is for the establishment of secure communication channels and TLS handshakes to ensure the integrity of the data exchanged. Afterwards, when the secure channel is established, packets with TLS protocols with application data are sent, which most likely include sensitive data related to pairing, including authentication credentials. This process mainly used TLSv1.3, with a couple of packets using TLSv1.2. 

The traffic with TLS traffic also included some Client Hello messages that reveal the domain name being queried, including graph.facebook.com, therefore are considered as partially readable, as seen in Figure \ref{fig:NP_bar_log_third_party_protocol_traffic}. However, in the figure, there is a similar amount of TCP, TLSv1.2, and TLSv1.3 traffic, most of which remains unreadable. Most of the third-party traffic was unreadable, as seen in the figure. The third-party traffic was mainly to Google, while none was to Spotify. However, the only packets captured for the packets to Google were TCP protocols used for connectivity checks and roughly 10 packets of secure communication. This suggests that communication is more likely part of general internet connectivity rather than a dependency on the pairing process for services like authentication. The content of the readable HTTP traffic was the same as in voice-based commands, but here, four packets of HTTP traffic were captured, meaning two sets of the same HTTP GET request and HTTP response packets. The sequence of packets, as well as the domain names revealed in the pairing process, would allow external observers to deduce that it is from the initial pairing of the devices and then monitor the pairing process or further interactions, which allows for better tracking and user profiling.



\section{Idle Traffic Analysis}

\begin{figure}[b]
    \centering
    \includegraphics[width=0.9\textwidth]{Figures/NI_bar_log_different_party_traffic.pdf}
    \caption{Public network traffic captured during the idle traffic experiment, with the data categorized by party type.}
    \label{fig:NI_bar_log_different_party_traffic}
\end{figure}

An 8-hour capture was conducted while the glasses and phone were idle, aiming to identify any background activity that could pose potential privacy concerns. Public and private network traffic was once again quite balanced, with public traffic having roughly 48\% and about 0.83 MB of traffic.

Figure \ref{fig:NI_bar_log_different_party_traffic} shows that nearly all of the traffic captured occurred in the first four hours of the capture. Compared to active traffic, idle traffic had a much higher percentage of third-party traffic, contributing slightly above 10\% of the total traffic volume. There is minimal support traffic, and they appeared only during the initial half-hour of the experiment.


Figure \ref{fig:NI_line_all_ip_address_traffic} illustrates the traffic volume for all addresses that contributed more than 1100 bytes of volume. As depicted in the figure, most of the fluctuations arose from fluctuations in first-party traffic. The highest traffic volume appeared at the start of the capture, coinciding with when the glasses were about to be closed, placed into the charging case, and the phone switched off. Thereafter, the traffic sent to Meta servers occurs in bursts rather than in a consistent flow, although the traffic volume dropped to nearly zero after the four-hour mark. The traffic bursts may represent updates being sent or synchronizations with the glasses. 

\begin{figure}[H]
    \centering
    \includegraphics[width=\textwidth]{Figures/NI_line_all_ip_address_traffic.pdf}
    \caption{Public network traffic that contributed more than 1100 bytes captured during the idle traffic experiment, with the data categorized by IP addresses. The legend displays each IP address along with its corresponding ISP or Service if known, and the Total Number of Bytes.}
    \label{fig:NI_line_all_ip_address_traffic}
\end{figure}

Another possibility is that an app is running some unknown and possibly unnecessary background task. Due to the traffic being encrypted, the exact content being sent is unknown. Although not armed with knowledge of the content of the packets, the traffic metadata still reveals the user's geological location, as nearly all traffic is routed to Zürich, Switzerland. Although it does not pinpoint the user's exact location, it still provides a general location to determine where the user is located. The red line, shown to have the most fluctuations, also exists during active captures, which are usually more active in minute 0.5, which is when powering on the glasses and opening the Meta View app. Although the traffic volume is not as high, this address continues to have a constant presence throughout the remaining time of the active captures. This suggests that this destination address could mainly be for connection purposes. The destination for this IP address is the Meta server in Zürich, Switzerland, suggesting that all Meta traffic routed to the US is due to active interactions with the glasses.

\begin{figure}[H]
    \centering
    \includegraphics[width=\textwidth]{Figures/NI_line_combined_third_party_traffic.pdf}
    \caption{Third-party traffic that contributed more than 1000 bytes captured during the idle traffic experiment, with the data categorized by IP addresses. The legend displays each IP address along with its corresponding ISP or Service if known, and the Total Number of Bytes.}
    \label{fig:NI_line_combined_third_party_traffic}
\end{figure}

Third-party addresses that contributed more than 1000 bytes during the idle traffic experiments are illustrated in Figure \ref{fig:NI_line_combined_third_party_traffic}. An intriguing note in the figure is that every single spike in the traffic is from a Google IP address, but each spike represents a different IP address. For idle traffic, the third-party traffic was nearly all unreadable, which was also observed for each spike, which were packets with protected payload with the QUIC protocol. The timing for the third-party traffic also aligns well with the peaks observed in first-party traffic, which implies interdependencies between first parties and third parties. As mentioned, a Firebase database was identified in the app analysis \ref{chap:appAnalysis}. It could be analytic events via Firebase that are triggered somehow, which results in these simultaneous first-party and third-party traffic occurrences. The locations to these addresses showed a similar pattern as in active traffic, with Google and, this time, Microsoft in the US, whilst the rest of the traffic was mainly routed to Zürich, Switzerland.

Finally, Figure \ref{fig:NI_bar_log_all_protocol_traffic} shows that the readability of all traffic captured during the idle traffic experiment was predominantly unreadable, with few packets being partially readable due to the exposed domain names, such as www.googleapis.com, in the payload. Unlike the active interaction experiments, there was no readable traffic. Another difference found was that idle traffic also captured non-readable DHCP traffic, which did not exist in the active traffic captures. However, these are DHCP ACK packets, which contain negligible information and serve only to confirm IP address assignment. The unreadable traffic and what can be understood from partially readable traffic have the same observations as in active traffic interactions.

\begin{figure}[H]
    \centering
    \includegraphics[width=0.9\textwidth]{Figures/NI_bar_log_all_protocol_traffic.pdf}
    \caption{All network traffic captured during the idle traffic experiment, with the data categorized by readability and protocol.}
    \label{fig:NI_bar_log_all_protocol_traffic}
\end{figure}



\section{Summary}
Nearly all of the captured traffic was encrypted and, therefore, unreadable, which made the payload content inaccessible. However, findings from the experiments demonstrated that, despite the content of the packet being unreadable and without the ability to decrypt the payloads, there are still risks and concerns for user and bystander privacy. The analysis showed that purely observing the traffic patterns and using the metadata from the traffic, including packet size, timing, domain names, and destinations of IP addresses, is enough to expose application usage, user behavior, and the geological location of the user, allowing for user tracking and profiling. For the Ray-Ban Meta glasses, a few interactions stand out, including media transfer and receiving messages on WhatsApp.

The inability to know the traffic also poses a privacy concern to users who themselves are not aware of and therefore not able to give consent to what user data or activities may be tracked or collected. The traffic showed minimal interactions with third-party destinations, which assists in reducing the number of possible attack vectors and protecting user data. However, it is plausible that the service providers may process the traffic routed to Google or Spotify. Also, although the traffic captured did die down after four hours, the lack of transparency surrounding background tasks during the idle state raises questions about whether there may be continuous data capture without the user's awareness, which could potentially impact bystander privacy if audio or images are inadvertently transmitted.
